\section{Implementation}
\label{sec:impl}

\paragraph{Library and integration.} The implementation uses a shared Java
\texttt{SocketChannel}~\cite{android_socketchannel} HTTP/1.1 sender and a TLS~1.2 adapter built on \texttt{SSLEngine}~\cite{java_sslengine}. The latter verifies the peer certificate and host, encrypts at most 4\,KiB of plaintext per wrap, and rechecks policy before each at-most-4\,KiB ciphertext write. Proxies, redirects, pooled writers, early data and renegotiation are excluded. Its positive write return is socket acceptance, not server delivery. A local
mutex plus file lock coordinates threads and processes. The ACRA and Matomo embeddings and their sizes are in
Appendix~\ref{app:sdk}.

\paragraph{Fixture app (\lab).} To compare enforcement points on identical schedules, we built a fixture app with ten
\emph{tenants}. Each is an independent workload with its own scope, stated contract and persisted jobs (Table~\ref{tab:tenants}).
Every tenant collects 12 records tagged with synthetic tracer strings such as \texttt{RD-pen007-T1-e0-03} into a SQLite
outbox. Tenants with a summary also derive one daily summary row whose content lists its source tags, and copy it into a
persisted summary job. All uploads use persisted \texttt{JobScheduler} jobs with a network constraint, in batches of four
records with up to five in-job attempts and exponential backoff. They stop cooperatively between batches on
\texttt{onStopJob}. The app's startup path re-enqueues an upload whenever its outbox holds unsent rows and no job is
scheduled, as analytics SDKs commonly do. One package is installed per enforcement variant, all from the same code:
\begin{itemize}[leftmargin=*,nosep]
\item \vname{none}: stops collection only.
\item \vname{cancel}: additionally cancels the withdrawn scope's jobs and purges its unsent outbox.
\item \vname{enqueue}: checks consent at every enqueue, including the startup path.
\item \vname{jobstart}: checks consent at job start.
\item \vname{conv}: cancel + enqueue + job-start checks, the strongest conventional combination.
\item \vname{egress}: a serialized egress check on the job's declared scope, without lineage or epochs.
\item \vname{rd}: \sys.
\end{itemize}
Two ablations change one \sys mechanism each. \vname{volatile} keeps labels in memory and fails open on unknown
provenance. \vname{taint} uses durable lineage but a contract-unaware rule, so any withdrawal or revocation of a source
blocks all its uses, as a deliberately over-broad synthetic policy would. It is not an implementation of TaintDroid or a comparison with a complete taint system. A tenth package, \vname{noser}, drops serialization and is
used only in the race micro-benchmark.

Five matched variants share durable lineage, epochs, and the same chunk transport:
\vname{menqueue}, \vname{mstart}, \vname{mlayer} (cancellation plus enqueue, job-start and retry checks),
\vname{mfence}, and \vname{mboth} (layered checks plus final fencing). All retain source labels
in copied job inputs. Their comparisons isolate final placement more closely than the boolean-only baselines.

\paragraph{Harness and oracle.} A host sink, reached via \texttt{10.0.2.2}, logs every request and waits 300\,ms before
responding, emulating a wide-area server. The oracle (\S\ref{sec:setup}) is a separate Python reimplementation of the
contract table. It never reads the app's labels or decisions, but takes order from device-logged events; the receiver is
the only outside observer. The split policy for mixed requests, the real-app embedding,
and the phone and throttled-link harnesses are described in Appendix~\ref{app:round3impl}.

\paragraph{Static scanner.} For the measurement (\S\ref{sec:measure}) we wrote a standard-library-only DEX parser. It
reads string, type, method and class tables, decodes Dalvik code items, and builds an over-approximate call graph:
class-hierarchy resolution below app-defined types, and allocation edges to callback methods of allocated classes.
